%============================================================================
\section{Conclusiones}\label{sec:conclusiones}
%============================================================================

\begin{enumerate}[label=\arabic*), leftmargin=*]
  \item El enlace BPSK con filtro acoplado toleró impulsos de casi el doble de
        la amplitud del símbolo sin errores: el umbral observado coincidió con
        $A\,\mathrm{SPS}/W=2$, y la constelación de decisión confirmó la
        amplitud efectiva $A_iW/\mathrm{SPS}$.
  \item Por encima del umbral la BER se estabilizó en $1{,}1\times10^{-2}$, sin
        crecer al duplicar la amplitud: la degradación impulsiva es de umbral y
        su piso depende de la fracción y polaridad de los símbolos alcanzados,
        no de la relación señal a ruido gaussiano.
  \item El filtro acoplado diluye los impulsos más cortos que el símbolo, pero
        hace que todo impulso influya en alguna decisión.
  \item La señal temporal, la constelación de decisión, el ojo y la BER
        revelaron la perturbación; la PSD no, por lo que las medidas promediadas
        no bastan para diagnosticar ruido impulsivo escaso.
  \item El modelo reprodujo las mediciones al considerar la secuencia de prueba
        periódica, que debe tratarse como parte del diseño experimental.
  \item La prueba en hardware con 16QAM mostró que, en un enlace físico, la
        relación señal a ruido en el receptor determina si la constelación es
        reconocible, lo que refuerza la importancia de la separación entre
        símbolos considerada al elegir la modulación.
\end{enumerate}

\textbf{Trabajo futuro.} Sería ideal implementar en hardware el sistema
completo en el que se basó este proyecto: el enlace BPSK del escenario S1 con
su receptor de filtro acoplado, transmitido entre radios HackRF One, para
analizar experimentalmente el desempeño de la modulación seleccionada. Esa
implementación permitiría medir la curva de BER frente a $E_b/N_0$ con la
instrumentación de la fase de validación del curso, introducir perturbaciones
impulsivas reales o inyectadas y compararla con 16QAM bajo las mismas
condiciones. También se propone modelar llegadas aleatorias de impulsos
(Middleton o Bernoulli-gaussiano), estudiar impulsos más largos que un símbolo
y el pulso RRC, y evaluar el recorte o borrado de muestras y la codificación con
entrelazado.
